\chapter{Single Source of Truth i Offline Caching (Wykład 13)}

W poprzednich rozdziałach nauczyliśmy się pobierać dane z Internetu (Retrofit) oraz zapisywać je w lokalnej bazie danych (Room). Do tej pory traktowaliśmy te technologie rozłącznie. Nasze aplikacje działały w trybie \textit{Online-Only} - jeśli telefon tracił zasięg, użytkownik widział pusty ekran lub komunikat błędu.

W tym rozdziale połączymy te technologie, aby stworzyć aplikację typu \textbf{Offline-First}.

\section{Offline-First}

Urządzenia mobilne, ze swojej natury, działają w środowisku o zmiennej jakości połączenia sieciowego. Użytkownik może wejść do windy, przejechać przez tunel w metrze lub po prostu znaleźć się w strefie słabego zasięgu.

W prostych aplikacjach przepływ danych wygląda następująco:
\begin{center}
\texttt{UI} $\leftrightarrow$ \texttt{ViewModel} $\leftrightarrow$ \texttt{Retrofit (Internet)}
\end{center}

Wady tego rozwiązania są krytyczne dla User Experience (UX):
\begin{itemize}
    \item \textbf{Brak danych offline:} Bez internetu aplikacja jest bezużyteczna.
    \item \textbf{Opóźnienia (Latency):} Nawet przy dobrym zasięgu, każde wejście na ekran wymaga oczekiwania na odpowiedź serwera (ładowanie spinnera).
    \item \textbf{Utrata danych:} Jeśli użytkownik wypełni formularz i kliknie \textit{Wyślij} w momencie utraty zasięgu, dane przepadają.
\end{itemize}

Paradygmat Offline-First odwraca tę logikę. Zakłada on, że brak dostępu do Internetu nie jest błędem, lecz normalnym stanem pracy aplikacji. W tym modelu aplikacja działa \textbf{zawsze}, niezależnie od sieci. Wyświetla dane, które pobrała wcześniej (cache), a Internet służy jedynie do \textit{synchronizacji} (aktualizacji) tych danych w tle.

\textbf{Główne założenia:}
\begin{enumerate}
    \item Aplikacja uruchamia się i pokazuje treść natychmiast (korzystając z lokalnej bazy danych).
    \item Jeśli sieć jest dostępna, pobieramy nowsze dane i aktualizujemy bazę.
    \item Jeśli sieć jest niedostępna, użytkownik nadal może przeglądać stare dane, a nawet modyfikować stan (np. dodawać do ulubionych) - zmiany te zostaną wysłane, gdy połączenie wróci.
\end{enumerate}

Aby zrealizować ten cel, musimy gruntownie przebudować architekturę naszych Repozytoriów, implementując wzorzec \textbf{Single Source of Truth (SSoT)}.

\section{Wzorzec Single Source of Truth (SSoT)}

Wzorzec \textbf{Single Source of Truth (SSoT)} (Pojedyncze Źródło Prawdy) to fundament nowoczesnych aplikacji mobilnych. Interfejs użytkownika (UI) oraz ViewModel nigdy nie powinny komunikować się bezpośrednio z siecią (API). Jedynym zaufanym źródłem danych dla widoku jest \textbf{lokalna baza danych}.

Oznacza to, że nawet jeśli mamy świetne połączenie z Internetem, dane najpierw trafiają do bazy (Room), a dopiero stamtąd są wyświetlane na ekranie. W tym podejściu przepływ danych przestaje być liniowy (Pobierz $\rightarrow$ Wyświetl). Dzieli się on na dwa niezależne procesy:

\begin{enumerate}
    \item \textbf{Ścieżka Odczytu (Read Path):}
    \begin{itemize}
        \item Odpowiedzialność: Dostarczanie danych do UI.
        \item Kierunek: \texttt{Room $\rightarrow$ Repository $\rightarrow$ ViewModel $\rightarrow$ UI}.
        \item Mechanizm: Reaktywność (\texttt{Flow}). ViewModel subskrybuje zmiany w tabeli. Gdy tylko coś zmieni się w bazie, UI odświeża się \textit{automatycznie}.
    \end{itemize}
    
    \item \textbf{Ścieżka Synchronizacji (Write Path):}
    \begin{itemize}
        \item Odpowiedzialność: Aktualizacja danych w bazie.
        \item Kierunek: \texttt{Retrofit $\rightarrow$ Repository $\rightarrow$ Room}.
        \item Mechanizm: Funkcje zawieszające (\texttt{suspend}). Repozytorium pobiera JSON z sieci, mapuje go na Encje i zapisuje w bazie.
    \end{itemize}
\end{enumerate}

\begin{figure}[h]
    \centering
    \includegraphics[width=0.7\textwidth]{fig/pum2-w13-ssotscheme.png}
    \caption{Architektura SSoT: UI obserwuje bazę danych, a sieć służy tylko do aktualizowania tej bazy w tle.}
    \label{fig:Figura ssotscheme}
\end{figure}

Zobaczmy, jak zmienia się interfejs naszego repozytorium:

\begin{minted}[frame=single,framesep=10pt]{kotlin}
// Rozdzielamy obserwację od odświeżania
class UserRepository(
    private val api: UserApi,
    private val dao: UserDao
) {
    // 1. Ścieżka Odczytu: Zwraca Flow z bazy danych.
    // Działa zawsze, nawet bez internetu (zwraca puste dane lub stare dane).
    fun getUsersStream(): Flow<List<UserEntity>> {
        return dao.getAllUsers()
    }

    // 2. Ścieżka Synchronizacji: Pobiera z sieci i zapisuje do bazy.
    // Nie zwraca danych! Zwraca Unit lub błąd sieci.
    suspend fun refreshUsers() {
        val remoteUsers = api.getUsers()
        // Logika zapisu do bazy (np. usuń stare, wstaw nowe)
        dao.replaceUsers(remoteUsers.map { it.toEntity() })
    }
}
\end{minted}

Metoda \texttt{refreshUsers()} nie zwraca listy użytkowników. Jej zadaniem jest jedynie aktualizacja bazy danych. Gdy baza zostanie zaktualizowana, Room automatycznie wyemituje nową listę w strumieniu \texttt{getUsersStream()}, co spowoduje odświeżenie ekranu.

\section{Podstawowy Caching}

Najprostszym sposobem na zsynchronizowanie lokalnej bazy danych z serwerem jest potraktowanie bazy jako tymczasowego lustra (cache). Przy każdym odświeżeniu usuwamy starą zawartość tabeli i wstawiamy nową, pobraną z API.

Algorytm metody \texttt{refresh()}:
\begin{enumerate}
    \item Pobierz listę obiektów z API.
    \item Otwórz transakcję bazy danych.
    \item Usuń wszystkie rekordy z tabeli (\texttt{DELETE FROM ...}).
    \item Wstaw nowe rekordy pobrane w kroku 1.
\end{enumerate}

Dlaczego usuwamy wszystko? Ponieważ na serwerze mogły zostać usunięte pewne elementy. Jeśli byśmy tylko używali \texttt{Insert (OnConflict = REPLACE)}, usunięte elementy \textit{wisiałyby} w naszej lokalnej bazie w nieskończoność.

W Room operacja \textit{Wyczyść i Wstaw} powinna być atomowa (wykonana w jednej transakcji), aby użytkownik nie zobaczył przez ułamek sekundy pustego ekranu między usunięciem a wstawieniem.

\begin{minted}[frame=single,framesep=10pt]{kotlin}
// data/local/UserDao.kt

@Dao
interface UserDao {
    // 1. Ścieżka Odczytu (zawsze zwraca najnowszy stan tabeli)
    @Query("SELECT * FROM users")
    fun getAllUsers(): Flow<List<UserEntity>>

    // Operacje pomocnicze dla Ścieżki Zapisu
    @Insert(onConflict = OnConflictStrategy.REPLACE)
    suspend fun insertAll(users: List<UserEntity>)

    @Query("DELETE FROM users")
    suspend fun clearAll()

    // Transakcja: Gwarantuje, że usunięcie i wstawienie
    // zdarzy się "jednocześnie" dla obserwatorów Flow.
    @Transaction
    suspend fun replaceUsers(users: List<UserEntity>) {
        clearAll()
        insertAll(users)
    }
}
\end{minted}

Repozytorium łączy API z DAO. Tutaj mapujemy też obiekty sieciowe (DTO) na obiekty bazy danych (Entity).

\begin{minted}[frame=single,framesep=10pt]{kotlin}
// data/UserRepository.kt

class UserRepository @Inject constructor(
    private val api: UserApi,
    private val dao: UserDao
) {
    // Ścieżka Odczytu - po prostu przekazujemy Flow
    val users: Flow<List<UserEntity>> = dao.getAllUsers()

    // Ścieżka Zapisu
    suspend fun refreshUsers() {
        try {
            // 1. Pobierz z sieci
            val remoteUsers = api.getUsers()
            
            // 2. Zmapuj na encje
            val entities = remoteUsers.map { dto -> 
                UserEntity(
                    id = dto.id,
                    name = dto.name,
                    email = dto.email
                    // Uwaga: Nie mapujemy pól lokalnych (np. isFavorite),
                    // bo API ich nie zwraca!
                ) 
            }
            
            // 3. Zaktualizuj bazę (Source of Truth)
            dao.replaceUsers(entities)
            
        } catch (e: Exception) {
            // Błąd sieci? Trudno.
            // Nie rzucamy wyjątku wyżej (chyba że chcemy pokazać Snackbar).
            // Dzięki temu UI nadal wyświetla stare dane z 'users' Flow.
            throw e 
        }
    }
}
\end{minted}

ViewModel uruchamia odświeżanie przy starcie, ale dane pobiera ze strumienia.

\begin{minted}[frame=single,framesep=10pt]{kotlin}
@HiltViewModel
class UserViewModel @Inject constructor(
    private val repository: UserRepository
) : ViewModel() {

    // UI obserwuje to pole. Będzie ono zawierać dane z bazy.
    // Dzięki stateIn, Flow zamienia się w StateFlow (Hot Stream).
    val uiState = repository.users
        .stateIn(viewModelScope, SharingStarted.WhileSubscribed(5000), emptyList())

    init {
        refreshData()
    }

    fun refreshData() {
        viewModelScope.launch {
            try {
                repository.refreshUsers()
            } catch (e: Exception) {
                // Tu obsłuż błąd, np. wyślij zdarzenie "Brak sieci" do UI.
                // Ale lista użytkowników NIE zniknie z ekranu!
            }
        }
    }
}
\end{minted}

Powyższa strategia jest świetna dla danych \textit{tylko do odczytu} (np. lista wiadomości). Ma jednak jedną krytyczną wadę.

Wyobraźmy sobie, że dodajemy do encji pole \texttt{isFavorite: Boolean}, które użytkownik zmienia klikając \textit{serduszko} w aplikacji.
\begin{enumerate}
    \item Użytkownik klika \textit{Ulubione} przy użytkowniku Jan Kowalski. W bazie \texttt{isFavorite} zmienia się na \texttt{true}.
    \item Następuje odświeżenie danych z sieci (metoda \texttt{refreshUsers}).
    \item Metoda \texttt{clearAll()} usuwa wszystkich użytkowników (wraz z informacją o ulubionych).
    \item Metoda \texttt{insertAll()} wstawia świeże dane z serwera. Serwer nie wie nic o \textit{ulubionych}, więc pole to przyjmuje wartość domyślną (\texttt{false}).
\end{enumerate}

\textbf{Efekt:} Odświeżenie danych resetuje stan lokalny aplikacji. Aby temu zapobiec, musimy zastosować bardziej zaawansowane podejście.

\section{Zaawansowany Caching}

Aby rozwiązać problem znikających \textit{serduszek} przy odświeżaniu danych, musimy fizycznie odseparować dane pochodzące z serwera od danych generowanych przez użytkownika. Zastosujemy architekturę opartą na dwóch tabelach:
\begin{enumerate}
    \item \textbf{Tabela Cache (np. \texttt{movies\_cache}):} Zawiera surowe dane z API (tytuł, opis, plakat). Ta tabela jest \textit{tymczasowa} - strategia \texttt{refresh()} może ją do woli czyścić i nadpisywać.
    \item \textbf{Tabela Stanu (np. \texttt{favorite\_ids}):} Zawiera tylko identyfikatory elementów oznaczonych przez użytkownika. Ta tabela jest trwała - proces synchronizacji z siecią nigdy jej nie modyfikuje.
\end{enumerate}

Ostateczny obiekt domenowy (\texttt{Movie}), który trafi do UI, powstanie w wyniku połączenia tych dwóch tabel w czasie rzeczywistym. Potrzebujemy dwóch osobnych encji w bazie danych Room.

\begin{minted}[frame=single,framesep=10pt]{kotlin}
// data/local/MovieCacheEntity.kt
@Entity(tableName = "movies_cache")
data class MovieCacheEntity(
    @PrimaryKey val id: Int,
    val title: String,
    val overview: String
    // Tu NIE MA pola isFavorite!
)

// data/local/FavoriteIdEntity.kt
@Entity(tableName = "favorite_ids")
data class FavoriteIdEntity(
    @PrimaryKey val id: Int
    // Ta tabela przechowuje tylko ID ulubionych filmów
)

// Model Domenowy (to, co widzi UI)
data class Movie(
    val id: Int,
    val title: String,
    val overview: String,
    val isFavorite: Boolean // To pole obliczymy dynamicznie
)
\end{minted}

Najważniejszą częścią jest zapytanie w DAO. Musimy pobrać wszystkie filmy z cache i \textit{dokleić} do nich informację, czy ich ID znajduje się w tabeli ulubionych. Używamy do tego klauzuli \texttt{LEFT JOIN}.

\begin{minted}[frame=single,framesep=10pt]{kotlin}
// data/local/MovieDao.kt

@Dao
interface MovieDao {

    // Ścieżka Odczytu: Łączymy dwie tabele
    // (f.id IS NOT NULL) zwróci 1 (true) jeśli znaleziono ID w ulubionych,
    // lub 0 (false) jeśli nie znaleziono.
    @Query("""
        SELECT 
            m.id, 
            m.title, 
            m.overview,
            (f.id IS NOT NULL) as isFavorite
        FROM movies_cache AS m
        LEFT JOIN favorite_ids AS f ON m.id = f.id
    """)
    fun getMoviesStream(): Flow<List<Movie>>

    // Operacje na Cache (Sieć)
    @Query("DELETE FROM movies_cache")
    suspend fun clearCache()

    @Insert(onConflict = OnConflictStrategy.REPLACE)
    suspend fun insertCache(movies: List<MovieCacheEntity>)

    @Transaction
    suspend fun refreshCache(movies: List<MovieCacheEntity>) {
        clearCache()
        insertCache(movies)
    }

    // Operacje na Ulubionych (Użytkownik)
    @Insert(onConflict = OnConflictStrategy.IGNORE)
    suspend fun addFavorite(id: FavoriteIdEntity)

    @Query("DELETE FROM favorite_ids WHERE id = :id")
    suspend fun removeFavorite(id: Int)
}
\end{minted}

Prześledźmy ten sam scenariusz co w poprzedniej sekcji:

\begin{enumerate}
    \item Użytkownik klika serduszko przy filmie o ID=5.
    \item Aplikacja dodaje rekord \texttt{FavoriteIdEntity(5)} do tabeli \texttt{favorite\_ids}.
    \item Room automatycznie przelicza zapytanie \texttt{LEFT JOIN}. Dla filmu ID=5 warunek \texttt{f.id IS NOT NULL} staje się prawdą. UI odświeża się, pokazując czerwone serce.
    \item \textbf{Następuje odświeżenie z sieci.}
    \item Metoda \texttt{refreshCache()} czyści tabelę \texttt{movies\_cache} i wstawia nowe dane.
    \item Tabela \texttt{favorite\_ids} \textbf{nie jest ruszana}.
    \item Room ponownie przelicza \texttt{LEFT JOIN}. Nowy rekord filmu z cache łączy się z istniejącym rekordem w ulubionych.
    \item UI odświeża się z nowymi opisami filmów, ale serduszko przy ID=5 \textbf{pozostaje na miejscu}.
\end{enumerate}

To podejście jest odporne na błędy i stanowi podstawę aplikacji Offline-First. W tym modelu, jeśli serwer przestanie zwracać film o ID=5 (zostanie on usunięty z oferty), zniknie on również z listy w aplikacji (bo \texttt{FROM movies\_cache} go nie zwróci), nawet jeśli jego ID wciąż jest w tabeli ulubionych. Jest to zazwyczaj pożądane zachowanie (tzw. \textit{Orphaned Record} w ulubionych nie przeszkadza).

\section{Współbieżność (Race Conditions)}

Mamy system, który działa poprawnie w 99\% przypadków. Jednak w specyficznych warunkach (np. wolna sieć + niecierpliwy użytkownik) może dojść do uszkodzenia danych. Rozważmy następujący scenariusz wyścigu (Race Condition):

\begin{enumerate}
    \item Aplikacja rozpoczyna pobieranie danych z sieci (trwa to 2 sekundy).
    \item W międzyczasie użytkownik klika \textit{Dodaj do ulubionych} przy filmie X.
    \item Aplikacja zapisuje ID filmu X do tabeli \texttt{favorite\_ids}.
    \item Pobieranie z sieci kończy się sukcesem.
    \item Aplikacja czyści tabelę \texttt{movies\_cache} i wstawia nowe dane.
\end{enumerate}

W powyższym scenariuszu wszystko jest w porządku, ponieważ tabele są rozłączne. Ale co się stanie, jeśli logika biznesowa będzie bardziej złożona?

Załóżmy, że zamiast osobnej tabeli, trzymamy liczbę polubień bezpośrednio w tabeli z filmami (bo API zwraca globalną liczbę lajków, a my chcemy ją lokalnie inkrementować).

\begin{minted}[frame=single,framesep=10pt]{kotlin}
// Scenariusz Błędu:
// 1. Użytkownik klika \textit{Like} (Lokalnie: likes = 10 + 1 = 11)
// 2. Repozytorium pobiera dane z sieci (Serwer: likes = 10)
// 3. Repozytorium nadpisuje lokalne zmiany danymi z serwera (likes = 10)
\end{minted}

W efekcie lajk użytkownika zostaje \textit{połknięty} przez odświeżenie danych. Nawet przy dwóch tabelach istnieje ryzyko niespójności, jeśli operacje nie są atomowe.

\begin{minted}[frame=single,framesep=10pt]{kotlin}
suspend fun toggleFavorite(id: Int) {
    // 1. Sprawdź czy jest w ulubionych
    val isFavorite = dao.isFavorite(id) 
    
    // <--- W tym momencie następuje przełączenie wątku (context switch)
    // <--- Inny proces usuwa ten film z bazy
    
    // 2. Jeśli był, to usuń. Jeśli nie, to dodaj.
    if (isFavorite) dao.remove(id) else dao.add(id)
}
\end{minted}

Aby zabezpieczyć się przed takimi sytuacjami, musimy zsynchronizować dostęp do bazy danych, upewniając się, że w danym momencie tylko jedna operacja modyfikująca (Zapis z Sieci lub Zapis od Użytkownika) jest wykonywana.

\section{Rozwiązanie 1: Mutex (Mutual Exclusion)}

Aby zapobiec wyścigom (Race Conditions), musimy zagwarantować, że operacje zapisu do bazy danych będą wykonywane \textbf{sekwencyjnie}, a nie równolegle. W świecie Kotlin Coroutines służy do tego narzędzie o nazwie \texttt{Mutex} (od ang. \textit{Mutual Exclusion} - Wzajemne Wykluczanie).

Działa ono jak zamek w drzwiach: tylko jedna osoba (korutyna) może być w środku (w sekcji krytycznej). Inne muszą czekać w kolejce, aż zamek zostanie zwolniony. Klasa \texttt{Mutex} pochodzi z biblioteki \texttt{kotlinx.coroutines.sync}. Używamy metody \texttt{withLock \{ ... \}}, która automatycznie zamyka zamek na początku bloku i otwiera go na końcu (nawet jeśli wystąpi wyjątek). Zmodyfikujmy nasze repozytorium:

\begin{minted}[frame=single,framesep=10pt]{kotlin}
// data/MovieRepository.kt

class MovieRepository @Inject constructor(
    private val api: MovieApi,
    private val dao: MovieDao
) {
    // Obiekt Mutex musi być polem klasy (współdzielony przez wszystkie wywołania)
    private val mutex = Mutex()

    // Operacja 1: Odświeżanie z sieci
    suspend fun refreshMovies() {
        // KROK A: Pobieranie z sieci (Długie!)
        // Robimy to POZA blokadą. Nie chcemy blokować UI, gdy czekamy na serwer.
        val remoteMovies = api.getMovies()
        val entities = remoteMovies.map { it.toEntity() }

        // KROK B: Zapis do bazy (Krótkie!)
        // To jest nasza SEKCJA KRYTYCZNA.
        mutex.withLock {
            dao.refreshCache(entities)
        }
    }

    // Operacja 2: Zmiana statusu ulubione
    suspend fun toggleFavorite(movieId: Int) {
        // Cała logika "sprawdź i zapisz" musi być atomowa.
        mutex.withLock {
            val isFavorite = dao.isFavorite(movieId)
            if (isFavorite) {
                dao.removeFavorite(movieId)
            } else {
                dao.addFavorite(FavoriteIdEntity(movieId))
            }
        }
    }
}
\end{minted}

Wróćmy do scenariusza wyścigu:
\begin{enumerate}
    \item \textbf{Wątek A (Refresh):} Pobiera dane z sieci (trwa to 2 sekundy). Mutex jest otwarty.
    \item \textbf{Wątek B (User):} Klika \textit{Like}. Wywołuje \texttt{toggleFavorite}.
    \item \textbf{Wątek B:} Wchodzi w \texttt{mutex.withLock}. Zamyka zamek. Zapisuje ulubiony film do bazy. Otwiera zamek.
    \item \textbf{Wątek A:} Kończy pobieranie. Próbuje wejść w \texttt{mutex.withLock}.
    \item Jeśli Wątek B jeszcze trzyma zamek, Wątek A \textbf{zawiesza się} (suspend) i czeka.
    \item Gdy Wątek B zwolni zamek, Wątek A wchodzi do sekcji krytycznej i bezpiecznie aktualizuje cache.
\end{enumerate}

Dzięki temu, że operacja zapisu ulubionego filmu (\texttt{addFavorite}) i czyszczenia cache (\texttt{refreshCache}) nigdy nie dzieją się w tym samym czasie, baza danych pozostaje spójna.

Zwróćmy uwagę, że pobieranie danych z sieci (\texttt{api.getMovies()}) znajduje się \textbf{przed} blokiem \texttt{mutex.withLock}. Gdybyśmy objęli blokadą również zapytanie sieciowe, użytkownik nie mógłby kliknąć \textit{Lubię to} przez cały czas trwania synchronizacji (aplikacja wydawałaby się zacięta). Blokujemy tylko dostęp do zasobu współdzielonego (Bazy Danych).

Zalety metody:
\begin{itemize}
    \item Prostota implementacji.
    \item Czytelność kodu (widać gdzie jest sekcja krytyczna).
\end{itemize}

Wady metody:
\begin{itemize}
    \item Ryzyko \textbf{zakleszczenia (Deadlock)}, jeśli spróbujemy zagnieździć dwa Mutexy lub wejść w ten sam Mutex dwukrotnie w tym samym wątku (choć Mutex w Kotlinie nie jest reentrant).
    \item Przy bardzo dużej liczbie operacji, ciągłe blokowanie i odblokowywanie może stać się wąskim gardłem wydajności.
\end{itemize}

\section{Rozwiązanie 2: Wzorzec Aktor (Actor Pattern)}

Mutex rozwiązuje problem, ale działa na zasadzie twardej blokady. Jeśli odświeżanie trwa 5 sekund, a użytkownik kliknie \textit{Lubię to}, operacja ta musi czekać w zawieszeniu.

Alternatywnym podejściem jest \textbf{Wzorzec Aktor}. Zamiast pozwalać wielu wątkom na bezpośredni dostęp do bazy danych (i kazać im czekać w kolejce), tworzymy jednego \textit{Aktora} - dedykowaną korutynę, która jako jedyna ma prawo pisać do bazy. Reszta aplikacji nie modyfikuje danych bezpośrednio, lecz wysyła do Aktora wiadomości (Prośby/Akcje).

Repozytorium przestaje być tylko zbiorem metod. Staje się aktywnym komponentem, który posiada:
\begin{itemize}
    \item \textbf{Kanał (Channel):} Skrzynka pocztowa, do której wpadają zlecenia (np. \textit{Odśwież}, \textit{Polub}).
    \item \textbf{Pętlę Aktora:} Nieskończona pętla wewnątrz \texttt{init}, która pobiera komunikaty jeden po drugim i je przetwarza.
\end{itemize}

Dzięki temu operacje są naturalnie zsynchronizowane (Aktor robi tylko jedną rzecz naraz), ale wątek wysyłający wiadomość nie jest blokowany - wrzuca list do skrzynki i idzie dalej. Do implementacji użyjemy \texttt{Channel} (kanał) oraz \texttt{Sealed Interface} do zdefiniowania typów wiadomości.

\begin{minted}[frame=single,framesep=10pt]{kotlin}
// data/MovieRepository.kt

// 1. Definiujemy możliwe akcje (Intent)
sealed interface DataAction {
    data object Refresh : DataAction
    data class ToggleFavorite(val id: Int) : DataAction
}

@Singleton
class MovieRepository @Inject constructor(
    private val api: MovieApi,
    private val dao: MovieDao,
    // Potrzebujemy Scope'a, aby uruchomić aktora, który żyje tak długo jak aplikacja
    @ApplicationScope private val externalScope: CoroutineScope
) {
    // 2. Kanał (Kolejka) o nielimitowanej pojemności
    private val actionChannel = Channel<DataAction>(Channel.UNLIMITED)

    // Ścieżka Odczytu (bez zmian)
    val movies: Flow<List<Movie>> = dao.getMoviesStream()

    init {
        // 3. Uruchomienie Aktora (Loop)
        externalScope.launch {
            // consumeEach przetwarza komunikaty sekwencyjnie (FIFO)
            actionChannel.consumeEach { action ->
                processAction(action)
            }
        }
    }

    // Logika przetwarzania (prywatna, dostępna tylko dla Aktora)
    private suspend fun processAction(action: DataAction) {
        when (action) {
            is DataAction.Refresh -> {
                try {
                    val remoteMovies = api.getMovies() // To może trwać długo
                    dao.refreshCache(remoteMovies.map { it.toEntity() })
                } catch (e: Exception) {
                    // Log error
                }
            }
            is DataAction.ToggleFavorite -> {
                val isFavorite = dao.isFavorite(action.id)
                if (isFavorite) dao.removeFavorite(action.id)
                else dao.addFavorite(FavoriteIdEntity(action.id))
            }
        }
    }

    // 4. Publiczne API (Fire and Forget)
    // Metody nie są już 'suspend', bo wrzucenie do kanału jest szybkie.
    fun refresh() {
        actionChannel.trySend(DataAction.Refresh)
    }

    fun toggleFavorite(id: Int) {
        actionChannel.trySend(DataAction.ToggleFavorite(id))
    }
}
\end{minted}

Wróćmy do naszego scenariusza wyścigu:
\begin{enumerate}
    \item UI wywołuje \texttt{refresh()} $\rightarrow$ Akcja \texttt{Refresh} wpada do kanału.
    \item Aktor wyjmuje \texttt{Refresh} i zaczyna pobierać dane z sieci.
    \item W międzyczasie UI wywołuje \texttt{toggleFavorite(5)} $\rightarrow$ Akcja \texttt{Toggle} wpada do kanału i czeka.
    \item Aktor kończy pobieranie i zapisuje dane cache.
    \item Aktor kończy przetwarzanie \texttt{Refresh}. Pętla wraca na początek.
    \item Aktor wyjmuje z kanału czekające \texttt{ToggleFavorite(5)}.
    \item Aktor aktualizuje tabelę ulubionych.
\end{enumerate}

Mimo że akcja użytkownika musiała \textit{poczekać} na przetworzenie, została ona zakolejkowana i wykonana poprawnie po zakończeniu synchronizacji. Żadne dane nie zostały utracone.

Zalety metody:
\begin{itemize}
    \item Brak klasycznego blokowania (deadlocks są niemożliwe).
    \item Gwarancja kolejności wykonywania operacji (FIFO).
    \item  UI nie czeka na wykonanie operacji zapisu (Fire and Forget).
\end{itemize}

Wady metody:
\begin{itemize}
    \item Brak natychmiastowej informacji zwrotnej (metoda \texttt{toggleFavorite} zwraca \texttt{void}). UI musi polegać na obserwacji \texttt{Flow} z bazy danych, aby dowiedzieć się, czy operacja się udała. Jest to zgodne z filozofią SSoT, ale wymaga zmiany myślenia programisty.
\end{itemize}

\section{Kompletny przykład: Czysta Architektura i Actor Pattern}

Poniżej znajduje się kompletny kod źródłowy modułu, który łączy architekturę Offline-First z pełnym podziałem na warstwy (Clean Architecture). 

W tym podejściu \textbf{ViewModel} nie rozmawia bezpośrednio z \textbf{Repozytorium}. Pośrednikiem są \textbf{Przypadki Użycia (Use Cases)}, które enkapsulują logikę biznesową.

\subsection{Warstwa Domeny (Domain Layer)}

Jest to warstwa niezależna od frameworków Androida. Definiuje modele, interfejsy oraz logikę biznesową.

Model Domenowy:
\begin{minted}[frame=single,framesep=10pt,breaklines]{kotlin}
// Plik: domain/model/Movie.kt
package com.example.movieapp.domain.model

data class Movie(
    val id: Int,
    val title: String,
    val overview: String,
    val rating: Double,
    val isFavorite: Boolean
)
\end{minted}

Interfejs Repozytorium (Kontrakt):
\begin{minted}[frame=single,framesep=10pt,breaklines]{kotlin}
// Plik: domain/repository/MovieRepository.kt
package com.example.movieapp.domain.repository

interface MovieRepository {
    // Strumień danych (zawsze aktualny)
    fun getMoviesStream(): Flow<List<Movie>>
    
    // Akcje (Fire-and-forget lub suspend)
    suspend fun refresh()
    suspend fun toggleFavorite(id: Int)
}
\end{minted}

Przypadki Użycia (Use Cases):
Każda klasa realizuje jedno zadanie biznesowe.

\begin{minted}[frame=single,framesep=10pt,breaklines]{kotlin}
// Plik: domain/usecase/GetMoviesUseCase.kt
package com.example.movieapp.domain.usecase

class GetMoviesUseCase @Inject constructor(
    private val repository: MovieRepository
) {
    operator fun invoke(): Flow<List<Movie>> {
        return repository.getMoviesStream()
    }
}

// Plik: domain/usecase/RefreshMoviesUseCase.kt
package com.example.movieapp.domain.usecase

class RefreshMoviesUseCase @Inject constructor(
    private val repository: MovieRepository
) {
    suspend operator fun invoke() {
        repository.refresh()
    }
}

// Plik: domain/usecase/ToggleFavoriteUseCase.kt
package com.example.movieapp.domain.usecase

class ToggleFavoriteUseCase @Inject constructor(
    private val repository: MovieRepository
) {
    suspend operator fun invoke(id: Int) {
        repository.toggleFavorite(id)
    }
}
\end{minted}

\subsection{Warstwa Danych (Data Layer)}

Implementacja dostępu do danych. Tutaj znajduje się logika \textbf{Room}, \textbf{Retrofit} oraz \textbf{Aktor}.

Encje Bazy Danych:
\begin{minted}[frame=single,framesep=10pt,breaklines]{kotlin}
// Plik: data/local/MovieEntities.kt
package com.example.movieapp.data.local

@Entity(tableName = "movies_cache")
data class MovieCacheEntity(
    @PrimaryKey val id: Int,
    val title: String,
    val overview: String,
    val voteAverage: Double
)

@Entity(tableName = "favorite_ids")
data class FavoriteIdEntity(
    @PrimaryKey val id: Int
)
\end{minted}

DAO (Data Access Object):
\begin{minted}[frame=single,framesep=10pt,breaklines]{kotlin}
// Plik: data/local/MovieDao.kt
package com.example.movieapp.data.local

@Dao
interface MovieDao {
    // SSoT: Łączenie danych z dwóch tabel
    @Query("""
        SELECT 
            m.id, m.title, m.overview, m.voteAverage as rating,
            (f.id IS NOT NULL) as isFavorite
        FROM movies_cache AS m
        LEFT JOIN favorite_ids AS f ON m.id = f.id
    """)
    fun getMoviesStream(): Flow<List<Movie>>

    // Operacje atomowe dla Aktora
    @Transaction
    suspend fun refreshCache(movies: List<MovieCacheEntity>) {
        clearCache()
        insertCache(movies)
    }

    @Query("DELETE FROM movies_cache")
    suspend fun clearCache()

    @Insert(onConflict = OnConflictStrategy.REPLACE)
    suspend fun insertCache(movies: List<MovieCacheEntity>)

    // Ulubione
    @Insert(onConflict = OnConflictStrategy.IGNORE)
    suspend fun addFavorite(entity: FavoriteIdEntity)

    @Query("DELETE FROM favorite_ids WHERE id = :id")
    suspend fun removeFavorite(id: Int)

    @Query("SELECT EXISTS(SELECT 1 FROM favorite_ids WHERE id = :id)")
    suspend fun isFavorite(id: Int): Boolean
}
\end{minted}

API (Retrofit):
\begin{minted}[frame=single,framesep=10pt,breaklines]{kotlin}
// Plik: data/remote/TmdbApi.kt
package com.example.movieapp.data.remote

data class TmdbResponse(@SerializedName("results") val results: List<MovieDto>)

data class MovieDto(
    val id: Int,
    val title: String,
    val overview: String,
    @SerializedName("vote_average") val voteAverage: Double
)

interface TmdbApi {
    @GET("movie/popular")
    suspend fun getPopularMovies(
        @Query("api_key") apiKey: String,
        @Query("language") language: String = "pl-PL"
    ): TmdbResponse
}
\end{minted}

Repozytorium (Implementacja Aktora):
Klasa ta implementuje interfejs z domeny. Używa kanału do synchronizacji.

\begin{minted}[frame=single,framesep=10pt,breaklines]{kotlin}
// Plik: data/repository/MovieRepositoryImpl.kt
package com.example.movieapp.data.repository

// Wewnętrzne wiadomości aktora
private sealed interface RepoAction {
    data object Refresh : RepoAction
    data class ToggleFavorite(val id: Int) : RepoAction
}

@Singleton
class MovieRepositoryImpl @Inject constructor(
    private val api: TmdbApi,
    private val dao: MovieDao,
    private val externalScope: CoroutineScope, // Application Scope
    private val apiKey: String
) : MovieRepository {

    private val actorChannel = Channel<RepoAction>(Channel.UNLIMITED)

    override fun getMoviesStream(): Flow<List<Movie>> = dao.getMoviesStream()

    init {
        // Uruchomienie pętli przetwarzania
        externalScope.launch {
            actorChannel.consumeEach { action ->
                processAction(action)
            }
        }
    }

    private suspend fun processAction(action: RepoAction) {
        when (action) {
            is RepoAction.Refresh -> {
                try {
                    val response = api.getPopularMovies(apiKey)
                    val entities = response.results.map { dto ->
                        MovieCacheEntity(dto.id, dto.title, dto.overview, dto.voteAverage)
                    }
                    dao.refreshCache(entities)
                } catch (e: Exception) { /* Ignoruj błąd sieci */ }
            }
            is RepoAction.ToggleFavorite -> {
                val isFav = dao.isFavorite(action.id)
                if (isFav) dao.removeFavorite(action.id)
                else dao.addFavorite(FavoriteIdEntity(action.id))
            }
        }
    }

    // Metody interfejsu wrzucają zadania do kolejki
    override suspend fun refresh() {
        actorChannel.send(RepoAction.Refresh)
    }

    override suspend fun toggleFavorite(id: Int) {
        actorChannel.send(RepoAction.ToggleFavorite(id))
    }
}
\end{minted}

\subsection{Dependency Injection (Hilt)}

Moduł spinający implementację Repozytorium z jego interfejsem.

\begin{minted}[frame=single,framesep=10pt,breaklines]{kotlin}
// Plik: di/AppModule.kt
package com.example.movieapp.di

@Module
@InstallIn(SingletonComponent::class)
object AppModule {

    @Provides
    @Singleton
    fun provideDatabase(@ApplicationContext context: Context): MoviesDatabase {
        return Room.databaseBuilder(context, MoviesDatabase::class.java, "app.db").build()
    }

    @Provides
    fun provideDao(db: MoviesDatabase): MovieDao = db.movieDao()

    @Provides
    @Singleton
    fun provideApi(): TmdbApi {
        return Retrofit.Builder()
            .baseUrl("https://api.themoviedb.org/3/")
            .addConverterFactory(GsonConverterFactory.create())
            .build()
            .create(TmdbApi::class.java)
    }

    @Provides
    @Singleton
    fun provideAppScope(): CoroutineScope = 
      CoroutineScope(SupervisorJob() + Dispatchers.Default)

    @Provides
    fun provideApiKey(): String = "YOUR_API_KEY"

    // Wstrzyknięcie Implementacji jako Interfejsu Domenowego
    @Provides
    @Singleton
    fun provideMovieRepository(
        impl: MovieRepositoryImpl
    ): MovieRepository = impl
}
\end{minted}

Powyższa definicja \texttt{provideAppScope} wymaga komentarza. Wstrzykiwanie \texttt{CoroutineScope} powiązanego z procesem aplikacji (a nie z ekranem) jest krytyczne dla wzorca Offline-First z dwóch powodów: 
\begin{enumerate}
    \item \textbf{Niezależność od cyklu życia UI (Fire-and-Forget):}
    Standardowy \texttt{viewModelScope} jest anulowany w momencie wyjścia użytkownika z ekranu lub rotacji urządzenia. Gdybyśmy użyli go do uruchomienia operacji zapisu (np. \texttt{toggleFavorite}), a użytkownik zamknąłby aplikację w trakcie trwania tej operacji, proces zapisu zostałby przerwany w połowie. Mogłoby to doprowadzić do niespójności danych. \texttt{externalScope} gwarantuje, że rozpoczęta operacja zostanie dokończona w tle, nawet jeśli interfejs użytkownika już nie istnieje.

    \item \textbf{Izolacja błędów (SupervisorJob):}
    Domyślny \texttt{Job} w korutynach działa na zasadzie odpowiedzialności zbiorowej - jeśli jedna korutyna w zakresie rzuci wyjątek, cały zakres jest anulowany (\textit{fail-fast}). Dla Aktora obsługującego wiele niezależnych zadań jest to zachowanie niepożądane (błąd synchronizacji sieci nie powinien blokować możliwości dodawania filmów do ulubionych). Użycie \texttt{SupervisorJob} sprawia, że awaria jednej korutyny nie wpływa na pozostałe ani na sam zakres, co zapewnia stabilność pętli przetwarzającej komunikaty Aktora.
\end{enumerate}

Linia kodu:
\begin{minted}{kotlin}
CoroutineScope(SupervisorJob() + Dispatchers.Default)
\end{minted}
składa się z trzech istotnych elementów, które razem tworzą bezpieczne środowisko dla operacji w tle. Pierwszy został omówiony powyżej, przyjrzyjmy się pozostałmy dwóm:

\begin{enumerate}
    
    \item \textbf{\texttt{Dispatchers.Default}}: 
    Definiuje domyślną pulę wątków dla tego zakresu. \texttt{Default} jest zoptymalizowany pod kątem operacji obciążających procesor (CPU-intensive), co jest dobrym punktem wyjścia dla ogólnej logiki biznesowej (np. przetwarzania logiki Aktora). Należy pamiętać, że operacje wejścia/wyjścia (I/O), takie jak zapytania do Room czy Retrofita, i tak powinny wewnętrznie przełączać się na \texttt{Dispatchers.IO}.
    
    \item \textbf{Operator \texttt{+}} : 
    W Kotlinie \texttt{CoroutineContext} działa jak mapa (zbiór elementów). Operator plus służy do łączenia tych elementów. W tym przypadku tworzymy kontekst, który posiada zarówno specyficzne zachowanie błędów (\texttt{SupervisorJob}), jak i specyficzną strategię wątkowania (\texttt{Dispatchers.Default}).
\end{enumerate}

Podsumowując: Tworzymy przestrzeń życiową dla wątków, która jest odporna na pojedyncze awarie i działa na wydajnej puli wątków w tle.

\subsection{Warstwa Prezentacji (ViewModel)}

ViewModel korzysta wyłącznie z Use Caseów. Nie ma dostępu do Repozytorium ani Bazy Danych.

\begin{minted}[frame=single,framesep=10pt,breaklines]{kotlin}
// Plik: presentation/MoviesViewModel.kt
package com.example.movieapp.presentation

@HiltViewModel
class MoviesViewModel @Inject constructor(
    getMoviesUseCase: GetMoviesUseCase,
    private val refreshMoviesUseCase: RefreshMoviesUseCase,
    private val toggleFavoriteUseCase: ToggleFavoriteUseCase
) : ViewModel() {

    // Konwersja Flow na StateFlow (UI State)
    val uiState: StateFlow<List<Movie>> = getMoviesUseCase()
        .stateIn(
            scope = viewModelScope,
            started = SharingStarted.WhileSubscribed(5000),
            initialValue = emptyList()
        )

    init {
        // Pobranie danych przy starcie
        viewModelScope.launch {
            refreshMoviesUseCase()
        }
    }

    fun onFavoriteClick(movie: Movie) {
        viewModelScope.launch {
            toggleFavoriteUseCase(movie.id)
        }
    }
    
    fun onSwipeRefresh() {
        viewModelScope.launch {
            refreshMoviesUseCase()
        }
    }
}
\end{minted}

\subsection{Konfiguracja Aplikacji (Hilt \& Manifest)}

Aby wstrzykiwanie zależności zadziałało, musimy utworzyć klasę dziedziczącą po \texttt{Application} i opatrzyć ją adnotacją \texttt{@HiltAndroidApp}. To tutaj Hilt generuje graf zależności.

\begin{minted}[frame=single,framesep=10pt,breaklines]{kotlin}
// Plik: MovieApplication.kt
package com.example.movieapp

@HiltAndroidApp
class MovieApplication : Application()
\end{minted}

Nie zapomnij zarejestrować tej klasy w pliku \texttt{AndroidManifest.xml}. Bez tego aplikacja wyrzuci błąd przy starcie.

\begin{minted}[frame=single,framesep=10pt,breaklines]{xml}
<manifest xmlns:android="http://schemas.android.com/apk/res/android"
    package="com.example.movieapp">

    <uses-permission android:name="android.permission.INTERNET" />

    <application
        android:name=".MovieApplication" 
        android:allowBackup="true"
        android:icon="@mipmap/ic_launcher"
        android:label="@string/app_name"
        android:theme="@style/Theme.MovieApp">
        
        <activity
            android:name=".presentation.MainActivity"
            android:exported="true">
            <intent-filter>
                <action android:name="android.intent.action.MAIN" />
                <category android:name="android.intent.category.LAUNCHER" />
            </intent-filter>
        </activity>
    </application>
</manifest>
\end{minted}

\subsection{Warstwa Widoku (UI - Jetpack Compose)}

Ostatnim elementem jest \texttt{MainActivity}, która pełni rolę punktu wejścia (\texttt{@AndroidEntryPoint}). Interfejs użytkownika budujemy w oparciu o \textbf{Jetpack Compose}, obserwując \texttt{StateFlow} z ViewModelu. Widok jest całkowicie pasywny - jedynie wyświetla stan i przekazuje zdarzenia (kliknięcia) do ViewModelu.

\begin{minted}[frame=single,framesep=10pt,breaklines]{kotlin}
// Plik: presentation/MainActivity.kt
package com.example.movieapp.presentation

@AndroidEntryPoint
class MainActivity : ComponentActivity() {
    override fun onCreate(savedInstanceState: Bundle?) {
        super.onCreate(savedInstanceState)
        setContent {
            MaterialTheme {
                MoviesScreen()
            }
        }
    }
}

@OptIn(ExperimentalMaterial3Api::class)
@Composable
fun MoviesScreen(
    viewModel: MoviesViewModel = hiltViewModel()
) {
    // Obserwacja stanu UI (zamiana Flow na State)
    val movies by viewModel.uiState.collectAsState()

    Scaffold(
        topBar = {
            TopAppBar(
                title = { Text("Filmy Offline-First") },
                actions = {
                    IconButton(onClick = { viewModel.onSwipeRefresh() }) {
                        Icon(Icons.Default.Refresh, contentDescription = "Odśwież")
                    }
                }
            )
        }
    ) { padding ->
        LazyColumn(
            contentPadding = padding,
            modifier = Modifier.fillMaxSize()
        ) {
            items(movies, key = { it.id }) { movie ->
                MovieItem(
                    movie = movie,
                    onFavoriteClick = { viewModel.onFavoriteClick(movie) }
                )
            }
        }
    }
}

@Composable
fun MovieItem(
    movie: Movie,
    onFavoriteClick: () -> Unit
) {
    ListItem(
        headlineContent = { Text(movie.title) },
        supportingContent = { 
            Text(
                text = movie.overview, 
                maxLines = 2, 
            ) 
        },
        trailingContent = {
            IconButton(onClick = onFavoriteClick) {
                // Ikona zmienia się w zależności od pola isFavorite (SSoT)
                Icon(
                    imageVector = if (movie.isFavorite) Icons.Default.Favorite 
                                  else Icons.Default.FavoriteBorder,
                    contentDescription = null,
                    tint = if (movie.isFavorite) Color.Red else Color.Gray
                )
            }
        },
        leadingContent = {
            Surface(
                color = MaterialTheme.colorScheme.primaryContainer,
                shape = MaterialTheme.shapes.small,
                modifier = Modifier.size(40.dp)
            ) {
                Box(contentAlignment = Alignment.Center) {
                    Text(
                        text = movie.rating.toString(),
                    )
                }
            }
        },
        modifier = Modifier.clickable { /* Otwórz szczegóły */ }
    )
    Divider()
}
\end{minted}